% sec-05-legislation.tex - Section 5, Legislation.  A main section.  FINAL stage.
% Figures 14, 15 and 16 drawn; Tables 14, 15 and 16 populated from the three
% archives under new-trial-system/inputs.  Subsection order is duties, drafting
% history, downward effect, which is also the order the float numbering fixes.
% Senior-author pass: Table 14's Duty column loses 0.10 cm to its Actor column,
% and two closing subsections state why a sponsor drafted a statute at all and
% where the legislative work is weakest.
\section{Legislation}
\label{sec:legislation}

\subsection{Who owes what}

Five versions of one bill and four companion documents were deposited between
May 30 and June 17, 2026. The bill is H. R. 9510, the Verification Before
Generation in Physical AI Oncology Trials Act of 2026, and by version 5.0 it is
a financial data amendment to the Federal Food, Drug, and Cosmetic Act
\cite{hr9510billv5}. Its findings state the operative requirement and the reason
for it in one sentence each. The requirement: Congress finds ``that an automated
process must verify robot-patient interaction code before that code is generated
or executed in a Physical AI oncology clinical investigation'' \cite{hr9510billv5}.
The reason: ``In a disembodied pipeline an error yields a wrong number a human
can catch downstream; in an embodied surgical system the same error class yields
a wrong motion a human cannot catch in time. No Federal law requires the
verification of robot-patient interaction code to clear before that code is
generated or executed, because the order of operations is not regulated''
\cite{hr9510billv5, verifygenact}.

The companion bill supplies the vocabulary the requirement needs. It defines
verification as ``confirmation that code was built correctly: implementation
conforms to its specification'' and validation as ``confirmation that the right
code was built, that is, the result agrees with an independent reference of
reality'', both against ASME V\&V 40, and it defines the class the gate applies
to: an artificial intelligence system ``whose output materially influences a
clinical decision or a robot action that can affect patient safety''
\cite{vvuqbill2026, asmevv40}. Figure 14 assigns each of the eleven duties the
bill creates to the actor who owes it, writing each duty inside the party that
owes it rather than beside it, so the six lines that remain in the figure are
only the places where two parties have to interact; and Table 14 lists the
statutory sections behind those duties.

% ---------------------------------------------------------------------------
% FIGURE 14.  plantuml-type, class diagram.  UPDATE STAGE REDRAW.  The final
% stage drew this as a use case diagram: six stick actors in two edge columns
% with thirteen association lines reaching across a system boundary into a field
% of eleven ellipses.  That drawing had two defects a reader sees immediately.
% The stick glyph is an elementary mark and reads as one beside a statutory
% argument, and an association drawn to a stick figure terminates at the label
% node rather than at the glyph, so the lines did not appear to connect.
%
% A class diagram removes both defects at their source.  Each actor is a class:
% a name compartment carrying the party, and a member compartment carrying the
% duties that party owes, written as members with their duty numbers.  Ownership
% is then read inside the box rather than traced along a line, and the only
% lines left are the six places where two parties actually interact.  Every one
% of the six runs edge to edge between two class frames, three vertically inside
% a column and three horizontally inside a row, so no line crosses another and
% none crosses a class.
%
% Geometry.  Three columns of 4.40 cm at x = 0.30, 5.85, 11.40, gutter 1.15 cm,
% which is the width a two-line association label needs plus 0.2 cm of clear
% canvas either side.  Two rows with their north west corners at y = 0 and
% y = -3.85, so the vertical associations have a 1.7 cm corridor for their
% labels.  Canvas 15.8 by 8.3 cm.
% ---------------------------------------------------------------------------
\begin{appfloat}[!tb]
\begin{appfig}[plantuml-type / class diagram]
\umlclass{sp}{0.30}{0}{4.40}{umlclshdr}{umlclsbody}{Sponsor}{%
  \guillemotleft regulated party\guillemotright\\[1.4pt]
  + U1 submit robot patient code\\
  + U3 publish the run cost ledger\\
  + U4 attest to VVUQ completion\\[1.4pt]
  prior law: partial U1, none U3, U4}
\umlclass{ci}{5.85}{0}{4.40}{umlclshdr}{umlclsbody}{Clinical investigator}{%
  \guillemotleft regulated party\guillemotright\\[1.4pt]
  + U1 submit robot patient code\\
  + U6 suspend on a failed gate\\[1.4pt]
  prior law: partial U1, none U6}
\umlclass{ib}{11.40}{0}{4.40}{umlclshdrg}{umlclsbodyg}{Institutional review board}{%
  \guillemotleft oversight body\guillemotright\\[1.4pt]
  + U5 review the Physical AI opt out\\
  + U6 suspend on a failed gate\\[1.4pt]
  prior law: partial U5, none U6}
\umlclass{se}{0.30}{-3.85}{4.40}{umlclshdrg}{umlclsbodyg}{Secretary, FDA}{%
  \guillemotleft Federal authority\guillemotright\\[1.4pt]
  + U7 maintain the national registry\\
  + U11 act on a financial deviation\\[1.4pt]
  prior law: none for either duty}
\umlclass{vs}{5.85}{-3.85}{4.40}{umlclshdrk}{umlclsbodyk}{Verification service}{%
  \guillemotleft actor the prior system does not define\guillemotright\\[1.4pt]
  + U2 clear verification before generation\\
  + U8 certify the test suite and coverage\\[1.4pt]
  prior law: no counterpart for either}
\umlclass{pt}{11.40}{-3.85}{4.40}{umlclshdr}{umlclsbody}{Participant}{%
  \guillemotleft rights holder\guillemotright\\[1.4pt]
  + U9 receive the autonomy disclosure\\
  + U10 exercise the opt out, no penalty\\[1.4pt]
  prior law: partial U9, none U10}
\begin{scope}[on background layer]
\node[umlpkg,fit=(sph)(sp)(cih)(ci)(ibh)(ib)(seh)(se)(vsh)(vs)(pth)(pt),
  inner sep=9pt] (bnd) {};
\end{scope}
\node[umlpkgtab,anchor=south west] at ([shift={(0mm,1.2mm)}]bnd.north west)
  {Verification Before Generation in Physical AI Oncology Trials Act, H. R. 9510};
% Three horizontal associations, each inside one row.  The label is set on two
% lines so that it clears both class frames inside the 1.15 cm gutter.
\draw[umlassoc] (sp.east) -- node[umlassoclbl,align=center,pos=0.5] {shares\\U1} (ci.west);
\draw[umlassoc] (ci.east) -- node[umlassoclbl,align=center,pos=0.5] {shares\\U6} (ib.west);
\draw[umlassoc] (se.east) -- node[umlassoclbl,align=center,pos=0.5] {registers\\U8} (vs.west);
% Three vertical associations, each inside one column.
\draw[umlassoc] (sp.south) -- node[umlassoclbl,pos=0.5] {U3 reported under U11} (seh.north);
\draw[umlassoc] (ci.south) -- node[umlassoclbl,pos=0.5] {U1 code enters the U2 gate} (vsh.north);
\draw[umlassoc] (ib.south) -- node[umlassoclbl,pos=0.5] {U5 protects U9 and U10} (pth.north);
% Annotation band, then the legend beneath it.  The note's leader is a short
% vertical run from the boundary edge to the note's own corner, so it crosses
% neither the legend row nor the in-figure note.
\node[umlnote,text width=50mm,anchor=north west] (n1) at (10.20,-6.80)
  {U2 and U8 are the two duties of the caption: no party in force today is
  positioned to discharge either one.};
\draw[trialslate,line width=0.4pt,dashed] (n1.north west) -- (vs.south east);
\pnote{0.30}{-7.15}{Ownership is read inside a class rather than along a line, so
only six associations are drawn\\where eleven duties are assigned. Three run
inside a row and three inside a column;\\none crosses another and none crosses a
class.}
\legkey{0.30}{-8.05}{trialburg}{regulated party or rights holder}
\legkey{4.60}{-8.05}{trialslate}{oversight body or Federal authority}
\legkey{9.20}{-8.05}{trialburgl}{party the prior trial system does not define}
\end{appfig}
\vspace{-0.6cm}
\figcaption{Figure 14. Six actors and eleven statutory duties, with the two
duties that\\no actor recognized by the prior trial system is positioned to
discharge.}
\end{appfloat}

\begin{apptable}
\begin{tabularx}{\textwidth}{>{\raggedright\arraybackslash}p{2.10cm} >{\raggedright\arraybackslash}p{4.85cm} >{\raggedright\arraybackslash}p{2.90cm} >{\raggedright\arraybackslash}X}
\headrow \thead{Section} & \thead{Duty created} & \thead{Actor} & \thead{Prior-law counterpart}\\
\midrule
Sec. 2 & Findings, the evidence to law to cost record & Congress & None \\
Sec. 3 & Amendment to the FD and C Act & Secretary & 21 U.S.C. as amended \\
Sec. 4 & Comparative provisions across prior versions & Secretary & Partial \\
Sec. 5 & Financial data, cost ledger per verification run & Sponsor & None \\
Appendix 6 & Cost estimate for the verification regime & Sponsor & None \\
Appendix 7 & Verification economics, per-run pricing & Sponsor & None \\
Appendix 8 & Financial standard for reported data & Secretary & Partial \\
Appendix 9 & Research matrix across prior developments & Secretary & None \\
Appendix 10 & Transparency of the reported ledger & Secretary & Partial \\
\end{tabularx}
\end{apptable}
\vspace{-0.6cm}
\tabcap{Table 14. The statutory sections, the duty each creates, the actor\\
who owes it, and whether prior law contains anything comparable.}

\subsection{How the text got there}

Figure 15 traces the drafting history, and Table 15 states the same lineage as
citable text. The eleven days between version 1.0 and version 5.0 are the
quantitative point: in the prior system, five substantive revisions of one bill
plus four companion documents is a legislative session's work rather than a
fortnight's. Each version's delta is traceable to a named companion. Version 4.0
put machine-readable diagrams inside statutory text, which the narrative case
document had argued for; version 5.0 attached a cost ledger to every
verification run, which the clinician trust framework and the enactment
framework had together made necessary
\cite{narrativecase2026, clintrust, congressvote}.

% ---------------------------------------------------------------------------
% FIGURE 15.  graphviz-type, cluster with records.  Five three-field version
% records at one width in an upper dashed cluster, four wider companion records
% below, one orthogonal supply edge routed 1.05 cm below the bill cluster.
% ---------------------------------------------------------------------------
\begin{appfloat}[!tb]
\begin{appfig}[graphviz-type / cluster with records]
\def\vw{2.70}
\foreach \i/\v/\d/\dl/\st in {%
  0/{v1.0}/{Jun 1 2026}/{standalone act}/{gvcells},
  1/{v2.0}/{Jun 4 2026}/{FD and C amendment}/{gvcells},
  2/{v3.0}/{Jun 4 2026}/{findings become figures}/{gvcellk},
  3/{v4.0}/{Jun 7 2026}/{Mermaid in statute}/{gvcellk},
  4/{v5.0}/{Jun 10 2026}/{cost ledger per run}/{gvcellh}}{
  \node[gvcellh,minimum width=\vw cm,minimum height=0.46cm,anchor=north west]
    at ({3.05*\i},0) {\v};
  \node[\st,minimum width=\vw cm,minimum height=0.46cm,anchor=north west]
    at ({3.05*\i},-0.46) {\d};
  \node[\st,minimum width=\vw cm,minimum height=0.46cm,anchor=north west]
    at ({3.05*\i},-0.92) {\dl};}
\begin{scope}[on background layer]
\node[gvcluster,minimum width=15.0cm,minimum height=1.90cm,anchor=north west] (bc) at (-0.15,0.15) {};
\end{scope}
\node[gvctitle,anchor=south west] at ([yshift=1.5mm]bc.north west) {Bill versions, H. R. 9510};
\foreach \i/\t/\d/\r in {%
  0/{VVUQ Physical AI Trial Bill}/{May 30 2026}/{statutory text and definitions},
  1/{From H. R. 9510 to Federal Law}/{Jun 14 2026}/{narrative case},
  2/{Earning the Clinician's Trust}/{Jun 16 2026}/{eight trust questions},
  3/{Earning the Congress's Vote}/{Jun 17 2026}/{enactment framework}}{
  \node[gvcellh,minimum width=3.2cm,minimum height=0.46cm,anchor=north west]
    at ({3.95*\i},-4.10) {\t};
  \node[gvcellg,minimum width=3.2cm,minimum height=0.46cm,anchor=north west]
    at ({3.95*\i},-4.56) {\d};
  \node[gvcellg,minimum width=3.2cm,minimum height=0.46cm,anchor=north west]
    at ({3.95*\i},-5.02) {\r};}
\begin{scope}[on background layer]
\node[gvcluster2,minimum width=15.0cm,minimum height=1.90cm,anchor=north west] (cc) at (-0.15,-3.95) {};
\end{scope}
\node[gvctitle,anchor=south west] at ([yshift=1.5mm]cc.north west) {Companion documents};
\foreach \i in {0,1,2,3}{
  \draw[gvedgeb] ({3.05*\i+\vw},-0.55) -- ({3.05*(\i+1)},-0.55);}
\draw[gvedged] (1.60,-4.10) -- node[left,font=\tiny,fill=trialwhite,inner sep=1pt] {supplies text} (1.35,-1.38);
\draw[gvedged] (5.55,-4.10) -- node[left,font=\tiny,fill=trialwhite,inner sep=1pt] {supplies case} (9.50,-1.38);
\draw[gvedged] (9.50,-4.10) -- (11.20,-2.95) -- node[right,font=\tiny,fill=trialwhite,inner sep=1pt] {supplies duties} (13.55,-1.38);
\draw[gvedged] (13.45,-4.10) -- node[right,font=\tiny,fill=trialwhite,inner sep=1pt] {supplies path} (13.90,-1.38);
\node[gvboxm,text width=54mm,anchor=north west] at (0,-6.35) {Five versions and four companions, May 30 to June 17, 2026};
\node[gvboxs,text width=44mm,anchor=north west] at (8.20,-6.35) {Eleven days between v1.0 and v5.0};
\pnote{0}{-7.20}{The five version records share one width and one three-field structure, so the
delta field reads across the rank\\as a single line. Companion records are 3.20 cm
wide against 2.70 cm, which separates the two classes without a second stroke.}
\end{appfig}
\vspace{-0.6cm}
\figcaption{Figure 15. Five bill versions across eleven days, the delta each
added, and\\the four companion documents that supplied the material for those
deltas.}
\end{appfloat}

\begin{apptable}
\begin{tabularx}{\textwidth}{>{\raggedright\arraybackslash}p{1.55cm} >{\raggedright\arraybackslash}p{2.15cm} >{\raggedright\arraybackslash}p{4.95cm} >{\raggedright\arraybackslash}X}
\headrow \thead{Version} & \thead{Deposit} & \thead{Delta over its predecessor} & \thead{Companion that supplied it}\\
\midrule
v1.0 & Jun 1, 2026 & Standalone Verification Before Generation Act & VVUQ Physical AI Oncology Trial Bill \\
v2.0 & Jun 4, 2026 & Recast as an FD and C Act amendment & Carried from v1.0 \\
v3.0 & Jun 4, 2026 & Findings rendered as figures in statutory text & Carried from v2.0 \\
v4.0 & Jun 7, 2026 & Machine-readable diagrams inside the bill & From H. R. 9510 to Federal Law \\
v5.0 & Jun 10, 2026 & Financial data amendment, cost ledger per run & Earning the Clinician's Trust; Earning the Congress's Vote \\
Total & Jun 1 to 10 & Five versions, four companions & Eleven days \\
\end{tabularx}
\end{apptable}
\vspace{-0.6cm}
\tabcap{Table 15. The five bill versions, the delta each added, and the companion\\document
that supplied the material, across eleven days in June 2026.}

\subsection{Where the requirement lands}

Figure 16 follows one requirement downward through four layers. The statute
names it, an adapted 21 CFR Subpart J constrains it, the Phase 0 gate of \S\ref{sec:protocol}
executes it, and a pre-incision checklist states it in the words a site operator
reads. The clinician trust framework supplies the bedside form of the same duty:
above the two supervision modes it defines, ``sits an unconditional control: the
clinician can override and stop the system at any time, for any reason, without
satisfying the gates first'' \cite{clintrust}. That sentence is what layer 4 of
Figure 16 is written to make executable, and it is why the legislative work in
this repository is not an abstraction: it terminates in a sentence somebody obeys
on the day of surgery. Table 16 places every legislative artifact on the same
nineteen-day window.

% ---------------------------------------------------------------------------
% FIGURE 16.  d2-type, layers.  Four equal bands of 1.85 cm at a uniform 0.75 cm
% gutter, three vertical descent edges at one x, no horizontal edges.
% ---------------------------------------------------------------------------
\begin{appfloat}[!tb]
\begin{appfig}[d2-type / layers]
\foreach \i/\ttl/\bx/\by/\bz/\st/\ts in {%
  0/{Layer 1, statute}/{H. R. 9510, Verification Before Generation Act}/{Amends the Federal Food, Drug, and Cosmetic Act}/{Financial data amendment, cost ledger per run}/{d2step}/{d2title},
  1/{Layer 2, regulation}/{21 CFR 312, Physical AI Subpart J}/{21 CFR 812 significant risk device pathway}/{Adapted ICH E6 R3 clinical practice text}/{d2ghost}/{d2title2},
  2/{Layer 3, protocol}/{Phase 0 simulation validation gate}/{At least 1000 simulations across 2 frameworks}/{Unified Safety Level at or above 7.0}/{d2ghost}/{d2title2},
  3/{Layer 4, site standard operating procedure}/{Pre-incision verification checklist}/{Do not generate motion code before the gate clears}/{Record the run cost and the verification hash}/{d2step}/{d2title}}{
  \begin{scope}[on background layer]
  \node[\st,minimum width=14.8cm,minimum height=1.85cm,anchor=north west]
    (bd\i) at (0,{-2.60*\i}) {};
  \end{scope}
  \node[\ts,anchor=north west] at ([shift={(2mm,-1.5mm)}]bd\i.north west) {\ttl};}
\foreach \i/\b/\st in {0/{H. R. 9510, Verification Before Generation Act}/{d2key},
  1/{21 CFR 312, Physical AI Subpart J}/{d2mid},
  2/{Phase 0 simulation validation gate}/{d2mid},
  3/{Pre-incision verification checklist}/{d2key}}{
  \node[\st,text width=42mm,anchor=north west] at (0.30,{-0.75-2.60*\i}) {\b};}
\foreach \i/\b in {0/{Amends the Federal Food, Drug, and Cosmetic Act},
  1/{21 CFR 812 significant risk device pathway},
  2/{At least 1000 simulations across 2 frameworks},
  3/{Do not generate motion code before the gate clears}}{
  \node[d2soft,text width=42mm,anchor=north west] at (5.15,{-0.75-2.60*\i}) {\b};}
\foreach \i/\b/\st in {0/{Financial data amendment, cost ledger per run}/{d2soft},
  1/{Adapted ICH E6 R3 clinical practice text}/{d2soft},
  2/{Unified Safety Level at or above 7.0}/{d2soft},
  3/{Record the run cost and the verification hash}/{d2gray2}}{
  \node[\st,text width=42mm,anchor=north west] at (10.00,{-0.75-2.60*\i}) {\b};}
\foreach \i/\lab in {0/{carries},1/{constrains},2/{executes}}{
  \draw[d2edgeb,line width=0.9pt] (2.40,{-1.85-2.60*\i}) --
    node[right,font=\tiny\sffamily,fill=trialwhite,inner sep=1pt] {\lab} (2.40,{-2.60-2.60*\i});}
\node[d2mid,text width=64mm,anchor=north west] at (5.15,-10.35)
  {Four equal bands, one uniform 0.75 cm gutter: a descent, not a funnel};
\pnote{0}{-11.05}{Layers 1 and 4 carry the solid burgundy stroke and layers 2 and 3 the ghost
stroke, because the statute and the\\operating room are where a requirement is
written and where it is obeyed; the two between are the machinery.}
\end{appfig}
\vspace{-0.6cm}
\figcaption{Figure 16. One requirement through four layers, from a statutory
sentence to\\the operating-room instruction, with the artifact that carries it at
each step.}
\end{appfloat}

\begin{apptable}
\begin{tabularx}{\textwidth}{>{\raggedright\arraybackslash}p{5.55cm} >{\raggedright\arraybackslash}p{1.95cm} >{\raggedright\arraybackslash}p{1.95cm} >{\raggedright\arraybackslash}X}
\headrow \thead{Legislative artifact} & \thead{Deposit} & \thead{Gap to next} & \thead{Prior-system counterpart}\\
\midrule
VVUQ Physical AI Oncology Trial Bill & May 30, 2026 & 2 days & A committee print cycle \\
Verification Before Generation Act of 2026 & Jun 1, 2026 & 3 days & A first draft \\
H. R. 9510 v3.0 & Jun 4, 2026 & 3 days & A markup revision \\
H. R. 9510 v4.0 & Jun 7, 2026 & 3 days & A markup revision \\
H. R. 9510 v5.0 & Jun 10, 2026 & 4 days & A reported bill \\
From H. R. 9510 to Federal Law & Jun 14, 2026 & 2 days & A committee report \\
Earning the Clinician's Trust & Jun 16, 2026 & 1 day & A stakeholder brief \\
Earning the Congress's Vote & Jun 17, 2026 & end & A floor memorandum \\
\end{tabularx}
\end{apptable}
\vspace{-0.6cm}
\tabcap{Table 16. Eight legislative artifacts in nineteen days, the gap between\\each
deposit, and the prior system's name for the equivalent stage.}

\subsection{What a bill has to survive}

A bill is not a paper, and the difference is worth stating because it governs
how the legislative artifacts in this repository were written. A paper is
finished when it is deposited. A bill is finished when it survives a markup, a
committee report, a cost estimate, and a floor vote, and each of those stages is
run by people who did not write it and who are entitled to ask what it costs.
The four companion documents exist because each answers one of those questions
before it is asked.

\emph{From H. R. 9510 to Federal Law} answers the narrative question: why this
statute rather than an agency guidance, and why a sequencing rule rather than a
performance standard \cite{narrativecase2026}. A performance standard says the
robot must be safe. A sequencing rule says the verification must clear before
the code is generated, which is checkable at a moment in time rather than
inferred from an outcome after the fact, and that is the property the embodiment
problem makes necessary.

\emph{Earning the Clinician's Trust} answers the adoption question through eight
framed questions, of which oversight and reliability are the two the statute
depends on \cite{clintrust}. Its oversight chapter distinguishes two supervision
modes: human-in-the-loop, where ``the system proposes an action under
verification before generation, but it cannot act until the clinician accepts
it'', and human-on-the-loop, where ``the system acts under audit within
pre-authorized bounds while the clinician monitors and retains the power to
intervene at any moment''. It then names the failure mode a statute cannot fix
by itself: over-trust, or automation bias, ``the tendency to defer to a machine
even when it is wrong'', against under-trust, ``the opposite error of discarding
useful guidance''. The design target is calibrated reliance, and the statutory
gate exists to make calibration possible rather than to replace it.

\emph{Earning the Congress's Vote} answers the enactment question, and the
financial appendices of version 5.0 answer the cost question directly: a
verification regime that cannot be priced per run cannot be scored, and a
provision that cannot be scored does not reach a floor \cite{congressvote,
hr9510billv5}. That is why the last substantive change to the bill in eleven
days of drafting was not a safety clause but a ledger clause. The three
appendices that carry it, cost estimate, verification economics, and financial
standard, are the ones a committee staff member reads first and the ones the
prior five versions did not have.

\subsection{What the legislative work does not claim}

Nothing in this repository is enacted, introduced, or endorsed. H. R. 9510 is a
drafting exercise carrying a plausible bill number, written to the form of a
House bill so that its provisions can be read against the Federal Food, Drug,
and Cosmetic Act as they would have to be read if they were ever offered
\cite{hr9510billv5}. The companion documents are analyses by a single author.
The value of the exercise is not that it will pass. It is that a sponsor
preparing a Physical AI oncology trial can now point at a specific statutory
sentence, a specific regulatory location, a specific protocol gate, and a
specific operating-room instruction, and show that the four agree; and a
regulator or a legislative staffer reading the same four can find the
disagreement, if there is one, in an afternoon rather than in a comment period.
That is the whole of the claim, and Figure 16 is the shape of it.

\subsection{Why a sponsor drafted a statute at all}

The obvious objection to \S\ref{sec:legislation} is that drafting legislation is
not a sponsor's job. It is not, and the reason it was done anyway is worth
stating because it generalizes. A sponsor preparing a Physical AI oncology trial
under existing law discovers that the thing it most needs to promise cannot be
promised in any recognized form. It can promise that its device meets a
performance specification, because 21 CFR \S 812 knows what that is. It can
promise that its drug is manufactured to a standard, because 21 CFR \S 312 knows
what that is. It cannot promise that the code governing robot-patient contact
was verified before it was generated, because no regulation names that ordering
and therefore no filing has a place to record it.

Faced with that gap a sponsor has three options. It can ignore the ordering and
promise only what the existing categories accept, which is what the prior system
does. It can invent a private assurance and ask a site to accept it, which
transfers the problem to an institutional review board with no basis for
evaluating it. Or it can write what the missing category would say, publish it,
and build its own protocol to the standard it proposes, so that the promise is
at least legible even where it is not yet required. The third is what this
repository did, and the Phase 0 gate of \S\ref{sec:protocol} is the same
requirement expressed in the one place a sponsor does control.

That is also why the bill's operative provision is a sequencing rule rather than
a performance standard. A performance standard would be satisfiable by a
sponsor's own testing after the fact, and would therefore add nothing to what 21
CFR \S 812 already asks. A sequencing rule is satisfiable only by evidence that
existed before the code did, which is a different kind of record and one a
verification service could be asked to hold. The two duties Figure 14 places in
a class no party in force today occupies are the whole content of that
difference.

\subsection{Where the legislative work is weakest}

Three weaknesses are worth conceding. The first is that a verification service
does not exist, and a statute that creates a duty no party is positioned to
discharge is not self-executing; the bill's appendices price the regime but do
not name who would be accredited to run it. The second is that the cost ledger
of version 5.0 assumes verification runs are priceable per run, which is true of
the author's own simulation record and may not be true of a service operating at
national scale. The third is that the eight trust questions of the clinician
framework were written by the same author as the bill they inform, so their
agreement is not independent evidence \cite{clintrust, hr9510billv5}.

None of the three is fatal to the exercise, because the exercise is not to enact
a statute. It is to give a sponsor, a site, and a reviewer one written object
that says what verification before generation would mean if it were required,
and to build a protocol that satisfies it whether or not it ever is. Table 16's
nineteen days measure how quickly that object can now be produced; the three
weaknesses measure how much of it would survive contact with a committee, and
the honest answer is that the sequencing rule would and the pricing might not.

\subsection{How the four documents divide the work}

The four legislative artifacts are not four drafts of one thing; each answers a
question the others cannot. Setting them side by side is the clearest way to see
what a sponsor-side legislative effort has to cover, and it is the reason
nineteen days produced four documents rather than one long one.

The VVUQ Physical AI Oncology Trial Bill supplies the definitions. Without it,
every operative sentence in H. R. 9510 would have to carry its own vocabulary,
and a bill that redefines verification, validation, uncertainty quantification,
high-risk artificial intelligence and AI-enabled device software function inside
its own text is a bill nobody can read against existing law. Anchoring those
terms to ASME V\&V 40 and to the device regulations means a reviewer can check
each definition against a standard rather than against the author's intent
\cite{vvuqbill2026, asmevv40}.

H. R. 9510 supplies the operative provision and, by version 5.0, the money. Its
findings section carries the engineering record the provision rests on, its
amendment section places the requirement inside the Federal Food, Drug, and
Cosmetic Act rather than beside it, and its financial appendices price the
regime per verification run \cite{hr9510billv5}. The placement matters as much
as the provision: a freestanding act creates a parallel regime that a sponsor
must satisfy twice, while an amendment creates one regime that the existing
inspection and enforcement machinery already knows how to reach.

\emph{From H. R. 9510 to Federal Law} supplies the narrative a committee needs,
and \emph{Earning the Clinician's Trust} supplies the bedside form of the duty,
which is where a statute either becomes practice or becomes paperwork
\cite{narrativecase2026, clintrust}. \emph{Earning the Congress's Vote} supplies
the enactment path \cite{congressvote}. Between them the five documents cover
definition, provision, placement, cost, narrative, practice and path, and Table
16 shows that the whole set was produced in the interval a single markup cycle
occupies in the prior system.

One consequence is worth drawing out for a funder rather than for a legislator.
A sponsor that can produce this set in nineteen days can also revise it in
nineteen days when a regulator, a site, or a committee objects. The prior
system's cost of legislative engagement is dominated by the cost of the next
draft, and a party that cannot afford the next draft argues its first one
longer than it should. That asymmetry is the same one \S\ref{sec:ind} identifies
in the documentation branch of the clinical hold fault tree, and it has the same
resolution: when re-issue is cheap, a defensible position is one you can revise
rather than one you have to defend.

